% =====================================================================================================
\section{Métodos implementados}\label{sec:metodos}

Esta seção descreve, um a um, os 29 métodos implementados. Começamos por um catálogo e por um exemplo numérico que
mostra por que famílias diferentes chegam a decisões diferentes; em seguida, detalhamos cada família.

\subsection{Interface comum}
Para que a comparação seja justa, todos os métodos usam a mesma interface. A cada rodada de decisão, recebem o mesmo
quadro de pistas e as mesmas zonas prioritárias, e devolvem uma ação --- uma pista ou a ação nula --- para cada
interceptador da fila. Há dois estilos de decisão:
\begin{itemize}
  \item Os métodos \conceito{individuais} (gulosos e árvores de comportamento) decidem um interceptador por vez, na ordem
        da fila, sem considerar os demais da rodada. A pista escolhida fica indisponível para os seguintes.
  \item Os métodos \conceito{conjuntos} resolvem a rodada inteira --- todos os interceptadores da fila contra todas as
        pistas --- e entregam a cada interceptador a sua parte da solução.
\end{itemize}
Nenhum método acessa a posição verdadeira das ameaças nem sabe quais pistas são falsas ou quais são iscas. Exceto M1 e M8, todos aplicam o
filtro de plausibilidade (E1).

\subsection{Catálogo}
A Tabela~\ref{tab:catalogo} lista os 29 métodos, numerados de M1 a M29, com a origem, o funcionamento e a adaptação de
cada um em uma linha. As subseções seguintes detalham cada um.

\begin{table}[p]
\centering
\caption{Catálogo dos métodos implementados. \emph{Ind.}: decide um interceptador por vez; \emph{conj.}: resolve a rodada
inteira. As referências são as de origem; a última coluna resume o que mudou para o nosso problema.}
\label{tab:catalogo}
\scriptsize
\renewcommand{\arraystretch}{1.25}
\begin{tabularx}{\textwidth}{@{}l>{\raggedright}p{2.6cm}>{\raggedright}p{1.6cm}>{\raggedright\arraybackslash}X>{\raggedright\arraybackslash}X@{}}
\toprule
\# & Método & Origem & Como funciona & Adaptação \\
\midrule
\multicolumn{5}{@{}l}{\emph{Regras simples}} \\
M1 & Aleatório & -- & pista sorteada, sem filtro, nunca troca (ind.) & piso de desempenho \\
M2 & Guloso (próximo) & \cite{kline2019} & menor $c_{ij}$ (ind.) & -- \\
M3 & Guloso (urgente) & \cite{kline2019} & menor folga até o solo (ind.) & folga estimada só pelo radar \\
M4 & Guloso (prioridade) & \cite{kline2019} & menor $c_{ij}/w_j$ (ind.) & peso inferido pela geometria (E3) \\
M5 & TEWA sequencial & \cite{roux2007} & ordena por ameaça, designa o mais rápido (conj.) & ameaça = peso / tempo até a zona \\
\midrule
\multicolumn{5}{@{}l}{\emph{Atribuição ótima (um interceptador por pista)}} \\
M6 & Húngaro & \cite{kuhn1955,munkres1957} & atribuição de mínimo tempo total (conj.) & custo com histerese (E2) \\
M7 & Húngaro ponderado & \cite{kuhn1955} & idem, tempo dividido pelo peso (conj.) & especialista das redes \\
M8 & Húngaro pond. sem filtro & \cite{kuhn1955} & M7 sem o filtro (E1) & mede o valor do filtro \\
M9 & Leilão & \cite{bertsekas1988} & lances e preços até o equilíbrio (conj.) & versão assimétrica sem escalonamento de $\varepsilon$ \\
\midrule
\multicolumn{5}{@{}l}{\emph{Alocação distribuída}} \\
M10 & Guloso global & \cite{choi2009} & fixa o par de maior lance, repete (conj.) & forma centralizada do CBAA \\
M11 & CBAA & \cite{choi2009} & lances locais + consenso de máximo (conj.) & comunicação completa via estação central \\
M12 & CBAA com comunic. local & \cite{choi2009} & idem, grafo de 3~km, 3 rodadas & conflitos resolvidos por ordem de chegada \\
\midrule
\multicolumn{5}{@{}l}{\emph{WTA (objetivo de valor esperado \eqref{eq:F})}} \\
M13 & MMR & \cite{denbroeder1959} & insere o par de maior ganho marginal (conj.) & custo de lançamento decide esperar \\
M14 & Húngaro sobre WTA & \cite{ahuja2007} & solução exata de \eqref{eq:F} com $\bar n=1$ (conj.) & coluna ``esperar'' por interceptador \\
M15 & MILP (HiGHS) & \cite{afonso2020,cruz2025} & mesmo problema como programa inteiro & parte de alocação do MPC inteiro \\
M16 & Busca local (VLSN) & \cite{ahuja2007} & melhora o MMR por realocação e troca & vizinhança polinomial \\
\midrule
\multicolumn{5}{@{}l}{\emph{Metaheurísticas (objetivo \eqref{eq:F}, semeadas com o MMR)}} \\
M17 & Algoritmo genético & \cite{lee2003ga} & torneio, cruzamento uniforme, mutação, elitismo & semente gulosa (``eugenia'') \\
M18 & Colônia de formigas & \cite{lee2002aco} & construção probabilística guiada por feromônio & opção ``esperar'' como caminho \\
M19 & Recozimento simulado & \cite{kline2019} & aceita pioras com probabilidade decrescente & parte do MMR \\
M20 & Enxame de partículas & \cite{kline2019} & partículas contínuas decodificadas por chave & codificação por chaves aleatórias \\
\midrule
\multicolumn{5}{@{}l}{\emph{Alocação dinâmica}} \\
M21 & Húngaro com reserva & \cite{hosein1990} & M7, mas guarda 20\,\% dos interceptadores & reserva só usada por pistas prioritárias \\
\midrule
\multicolumn{5}{@{}l}{\emph{Arquiteturas comportamentais}} \\
M22 & Árvore (distância) & \cite{ricardo2023} & engaja a mais próxima a menos de 4~km (ind.) & sem líder--seguidor \\
M23 & Árvore (utilidade) & \cite{feitosa2025} & avança se a utilidade afim $>0$ (ind.) & sem ramo de reforço; GA no treino \\
\midrule
\multicolumn{5}{@{}l}{\emph{Políticas aprendidas (imitam M7 e/ou otimizam o dano)}} \\
M24 & Imitação, DeepSets & \cite{ross2011,zaheer2017} & DAgger sobre M7 & amostragem estratificada das esperas \\
M25 & Imitação, atenção & \cite{ross2011,kool2019} & idem, codificador Transformer & idem \\
M26 & Imitação+PPO, DeepSets & \cite{schulman2017} & M24 refinada por reforço & perda auxiliar de legitimidade \\
M27 & Imitação+PPO, atenção & \cite{schulman2017,jia2026} & M25 refinada por reforço & idem \\
M28 & PPO, DeepSets & \cite{liu2023drones,palmas2025} & reforço do zero & política compartilhada \\
M29 & PPO, atenção & \cite{jia2026} & reforço do zero & política compartilhada \\
\bottomrule
\end{tabularx}
\end{table}

% ---------------------------------------------------------------------------------------------------------------
\subsection{Exemplo ilustrativo}\label{sec:exemplo}
Antes do detalhe, um exemplo pequeno mostra por que as famílias chegam a decisões diferentes. Há três interceptadores (A,
B e C) e três pistas; a pista 3 está numa zona prioritária, de modo que os pesos são $w=(1,1,10)$. Os tempos de
interceptação (em segundos) e as probabilidades de sucesso dadas por~\eqref{eq:q} são
\begin{equation*}
  c=\begin{pmatrix} 10 & 12 & 50\\ 11 & 60 & 70\\ 40 & 15 & 30\end{pmatrix},
  \qquad
  q=\begin{pmatrix} 0{,}836 & 0{,}833 & 0{,}782\\ 0{,}835 & 0{,}769 & 0{,}756\\ 0{,}795 & 0{,}829 & 0{,}809\end{pmatrix},
\end{equation*}
com uma linha por interceptador e uma coluna por pista. A Tabela~\ref{tab:exemplo} traz a solução de cada família,
obtida com as próprias implementações.

O guloso por proximidade, decidindo A primeiro, dá a A a pista 1 (10~s), embora A seja quase tão rápido para a pista 2
(12~s); com isso, B --- que só é rápido para a pista 1 --- fica com a pista 2, a 60~s. O húngaro enxerga a rodada inteira,
troca A e B de pista e reduz o tempo total de 100 para 53~s. O CBAA e o MMR também fixam primeiro o par mais atraente e
chegam à mesma solução que o guloso, embora por motivos diferentes: o MMR começa pela pista prioritária, cujo ganho
($10\times0{,}809$) domina os demais. O WTA exato encontra o maior valor esperado destruído (9,753 contra 9,691), o que
confirmamos enumerando as seis atribuições possíveis. O guloso por prioridade, ao dividir os tempos pelos pesos, manda o
interceptador mais rápido de todos (A) à pista prioritária e piora o restante.

\begin{table}[t]
\centering
\caption{Exemplo ilustrativo: solução de cada família para a mesma instância (pista 3 prioritária). $F$ é o valor
esperado destruído~\eqref{eq:F} sem custo de lançamento.}
\label{tab:exemplo}
\small
\begin{tabular}{@{}lcccl@{}}
\toprule
Método & Designação & Tempo total (s) & $F$ & Por quê \\
\midrule
Guloso próximo (ordem A,B,C) & A$\to$1, B$\to$2, C$\to$3 & 100 & 9,691 & A decide primeiro \\
Guloso prioridade & A$\to$3, B$\to$1, C$\to$2 & 76 & 9,484 & o mais rápido vai à prioritária \\
Guloso global / CBAA & A$\to$1, B$\to$2, C$\to$3 & 100 & 9,691 & fixa o melhor par primeiro \\
MMR & A$\to$1, B$\to$2, C$\to$3 & 100 & 9,691 & começa pela prioritária \\
Húngaro / leilão & A$\to$2, B$\to$1, C$\to$3 & 53 & 9,753 & ótimo do tempo total \\
WTA exato (M14, M15) & A$\to$2, B$\to$1, C$\to$3 & 53 & 9,753 & ótimo do valor esperado \\
\bottomrule
\end{tabular}
\end{table}

% ---------------------------------------------------------------------------------------------------------------
\subsection{Regras simples (M1--M5)}

\begin{metodo}{M1 -- Aleatório}
\ideia Medir o que se obtém sem usar informação alguma do quadro de radar.
\func Cada interceptador ocioso é lançado contra uma pista sorteada uniformemente entre as livres, sem filtro de
plausibilidade, e nunca troca de alvo.
\espera É o \emph{piso} de desempenho: um método que não o supera não está extraindo informação útil das pistas.
\end{metodo}

\begin{metodo}{M2--M4 -- Gulosos}
\ideia Cada interceptador escolhe, sozinho, a pista que lhe parece melhor por um critério simples.
\orig Regras de despacho comuns em sistemas reais~\cite{roux2007,kline2019}.
\func O interceptador da vez escolhe, entre as pistas livres e plausíveis,
\begin{equation}
  \text{M2: }\argmin_j c_{ij},\qquad
  \text{M3: }\argmin_j\,(t^G_j-c_{ij}),\qquad
  \text{M4: }\argmin_j \frac{c_{ij}}{w_j},
\end{equation}
em que $\argmin_j$ indica a pista $j$ que minimiza a expressão. M2 (\emph{próximo}) escolhe a pista que alcança mais
cedo. M3 (\emph{urgente}) escolhe a pista com a menor folga entre o tempo que ela leva até o solo,
$t^G_j=h_j^{\mathrm{AGL}}/(-v_j^z)$ (altura dividida pela velocidade de descida, infinito se a pista não desce), e o tempo de
interceptação: uma folga pequena significa que é preciso agir já. M4 (\emph{prioridade}) divide o tempo pelo peso da zona,
de modo que pistas prioritárias parecem dez vezes mais próximas. A pista escolhida fica indisponível para os
interceptadores seguintes da fila.
\prop Custo de $O(K)$ operações por interceptador, para $K$ pistas. Nenhuma garantia de qualidade; o resultado depende da
ordem em que os interceptadores decidem (Tabela~\ref{tab:exemplo}).
\adap O tempo até o solo (M3) e o peso (M4) são inferidos só pelo radar, porque o defensor não conhece o ponto de impacto.
\espera Rápidos e razoáveis quando há folga de interceptadores. M3 é o mais frágil diante de manobras, porque a velocidade
de descida estimada é ruidosa em zigue-zague e em voo baixo.
\end{metodo}

\begin{metodo}{M5 -- TEWA sequencial}
\ideia Primeiro decidir \emph{quais ameaças importam mais}; depois designar a cada uma o melhor interceptador disponível.
\orig Estrutura clássica dos sistemas de avaliação de ameaças e atribuição de armas~\cite{roux2007}.
\func Cada pista recebe uma nota de ameaça, que cresce quanto mais perto ela estiver de uma zona prioritária:
\begin{equation}
  \theta_j=\begin{cases} w_P/(1+t^Z_j), & t^Z_j<\infty,\\ w_O/10^3, & \text{caso contrário.}\end{cases}
\end{equation}
As pistas são visitadas em ordem decrescente de $\theta_j$, e cada uma recebe, entre os interceptadores ainda livres, o
de menor custo $c_{ij}$. O processo para quando acabam os interceptadores da rodada.
\prop Custo de $O(K\log K+nK)$ (ordenar as pistas e procurar o melhor interceptador para cada uma). Nenhuma garantia.
\adap A ameaça usa o tempo até a zona prioritária (E3) no lugar do tempo até o ativo defendido.
\espera Protege bem as zonas prioritárias, mas pode gastar interceptadores distantes com ameaças de nota alta enquanto
deixa passar ameaças de nota baixa e fáceis de pegar; herda os erros de $t^Z_j$ quando a rota passa perto de ativos
que não são o alvo.
\end{metodo}

% ---------------------------------------------------------------------------------------------------------------
\subsection{Atribuição ótima (M6--M9)}\label{sec:met_atrib}

\begin{metodo}{M6--M8 -- Método húngaro}
\ideia Escolher, para a rodada inteira, o conjunto de pares interceptador--pista de menor custo total, considerando todos
ao mesmo tempo.
\orig Kuhn~\cite{kuhn1955}, que o batizou em homenagem a trabalhos de matemáticos húngaros, e Munkres~\cite{munkres1957}.
\func O método resolve o problema de atribuição
\begin{equation}
  \min_{x}\sum_{i,j}\frac{c_{ij}}{w_j^{\,\beta}}x_{ij}\quad\text{s.a.}\quad
  \sum_j x_{ij}\le1,\ \ \sum_i x_{ij}\le1,\ \ x_{ij}\in\{0,1\},
  \label{eq:lap}
\end{equation}
isto é, minimiza a soma dos custos dos pares escolhidos, com cada interceptador e cada pista usados no máximo uma vez. O
problema pode ser \emph{retangular} (mais interceptadores que pistas, ou o contrário). A intuição do algoritmo é a
seguinte: subtrair uma constante de todos os custos de uma linha (ou de uma coluna) não muda qual atribuição é a melhor,
porque todas as atribuições usam exatamente uma posição daquela linha. O algoritmo subtrai constantes --- chamadas
\emph{potenciais} --- até que exista uma atribuição completa usando só posições de custo zero; como os custos restantes
são não negativos, essa atribuição é ótima. Encontrar a atribuição entre os zeros é feito por \emph{caminhos
alternantes}, que trocam pares já escolhidos por outros até acomodar mais um interceptador. Um resultado de teoria de
otimização (a \emph{unimodularidade total} da matriz de restrições) garante que, nesse problema, a solução ótima com
variáveis contínuas já é inteira, o que explica por que ele é fácil apesar de combinatório.
\prop Exato para~\eqref{eq:lap}, em $O(n^2K)$ operações para $n$ interceptadores e $K$ pistas.
\adap M6 usa $\beta=0$ (minimiza o tempo total); M7 usa $\beta=1$ (divide o tempo pelo peso, dando preferência às pistas
prioritárias); M8 é o M7 sem o filtro (E1), para medir quanto o filtro vale. A pista atual de um interceptador em voo entra
com a histerese (E2). M7 é também o especialista imitado pelas redes neurais. Usamos a implementação de
\texttt{scipy.optimize.linear\_sum\_assignment}.
\espera A solução é ótima para a rodada, mas não para o episódio: não considera ameaças que ainda vão aparecer, e o custo é
uma previsão de velocidade constante.
\end{metodo}

\begin{metodo}{M9 -- Leilão de Bertsekas}
\ideia Resolver o mesmo problema de atribuição como um mercado: os interceptadores ``compram'' as pistas, e os preços sobem
até que ninguém queira trocar.
\orig Bertsekas~\cite{bertsekas1988}.
\func Cada pista $j$ tem um preço $p_j$, inicialmente zero, e cada par tem um benefício $a_{ij}=-c_{ij}/w_j$ (quanto
menor o custo, maior o benefício). O \emph{valor líquido} de uma pista para um interceptador é o benefício menos o preço.
Um interceptador sem pista encontra a pista de maior valor líquido, $j_i$, com valor $v^{(1)}$, e a segunda melhor, com
valor $v^{(2)}$, e dá o lance
\begin{equation}
  p_{j_i}\leftarrow p_{j_i}+v^{(1)}-v^{(2)}+\varepsilon .
  \label{eq:lance}
\end{equation}
Ou seja, eleva o preço da pista até o ponto em que ela ainda lhe vale tanto quanto a segunda opção, mais um pequeno
incremento $\varepsilon$. Se a pista tinha dono, o dono anterior fica sem pista e volta a dar lances
(Algoritmo~\ref{alg:leilao}). O incremento $\varepsilon>0$ garante que os preços sempre sobem e que o processo termina.
\prop Ao final, cada interceptador está com uma pista cujo valor líquido está a menos de $\varepsilon$ do melhor
disponível para ele (condição chamada de \emph{$\varepsilon$-folga complementar}), e o custo total fica a menos de
$n\varepsilon$ do ótimo~\cite{bertsekas1988}. Com $\varepsilon$ pequeno, o leilão chega à mesma solução que o húngaro. Como
cada lance só usa os preços correntes, o algoritmo se presta a implementação paralela ou distribuída.
\adap Usamos os mesmos benefícios de M7, para que os dois sejam diretamente comparáveis. O caso \emph{assimétrico} (números
diferentes de interceptadores e pistas) exige cuidado. A prática usual de começar com $\varepsilon$ grande e reduzi-lo
(escalonamento de $\varepsilon$) acelera o leilão, mas, no caso assimétrico, deixa preços altos em pistas que terminam sem
dono e perde a garantia de otimalidade --- uma primeira versão com escalonamento coincidiu com o húngaro em só 10\,\% de
instâncias aleatórias. A versão final começa com preços nulos e usa um único $\varepsilon$ pequeno.
\parm $\varepsilon=10^{-3}\,\mathrm{amplitude}(a)/(n+1)$, em que a amplitude é a diferença entre o maior e o menor benefício.
\end{metodo}

\begin{algorithm}[H]
\caption{Leilão direto para atribuição assimétrica (M9).}\label{alg:leilao}
\begin{algorithmic}[1]
\Require benefícios $a_{ij}$ ($n\le m$), $\varepsilon>0$
\State $p_j\gets0$ para todo $j$; todos os $i$ livres
\While{existe $i$ livre com algum $a_{ij}$ finito}
  \State $j_i\gets\argmax_j(a_{ij}-p_j)$;\quad $v^{(1)}\gets a_{ij_i}-p_{j_i}$;\quad $v^{(2)}\gets\max_{j\ne j_i}(a_{ij}-p_j)$
  \State $p_{j_i}\gets p_{j_i}+v^{(1)}-v^{(2)}+\varepsilon$
  \State designa $j_i$ a $i$; se $j_i$ tinha dono $k$, $k$ volta a ficar livre
\EndWhile
\Ensure designação $i\mapsto j_i$
\end{algorithmic}
\end{algorithm}

% ---------------------------------------------------------------------------------------------------------------
\subsection{Alocação distribuída (M10--M12)}

\begin{metodo}{M10--M12 -- Guloso global e CBAA}
\ideia Permitir que cada interceptador escolha a sua pista trocando mensagens só com os vizinhos, sem decisor central, e
ainda assim evitar que dois interceptadores persigam a mesma pista.
\orig Choi, Brunet e How~\cite{choi2009}.
\func Cada interceptador $i$ atribui a cada pista um \emph{lance} $b_{ij}=w_j/\max(c_{ij},1)$, maior para pistas
prioritárias e próximas, e guarda três informações: a sua tarefa $x_i$, o maior lance que conhece para cada pista,
$y_{ij}$, e quem deu esse lance, $z_{ij}$. Duas fases se alternam (Algoritmo~\ref{alg:cbaa}):
\begin{itemize}
  \item na fase de \emph{lances}, um interceptador sem tarefa escolhe a pista em que o seu lance supera o maior lance que
        ele conhece, $x_i=\argmax_j\{b_{ij}:b_{ij}>y_{ij}\}$, e se declara vencedor dela;
  \item na fase de \emph{consenso}, cada interceptador compara o que sabe com o que sabem os vizinhos do grafo de
        comunicação e adota, para cada pista, o maior lance conhecido,
        \begin{equation}
          y_{ij}\leftarrow\max_{k\in\mathcal{N}_i\cup\{i\}}y_{kj},
          \label{eq:consenso}
        \end{equation}
        com desempate pelo menor índice. Se descobre que alguém deu um lance maior pela sua pista, desiste dela e volta a
        dar lances na rodada seguinte.
\end{itemize}
O processo para quando ninguém muda de tarefa.
\prop Como a informação sobre o maior lance se propaga um vizinho por rodada, o CBAA converge em no máximo (número de
agentes) $\times$ (diâmetro do grafo) rodadas, onde o diâmetro é o maior número de saltos entre dois agentes. O resultado
não tem conflitos e vale pelo menos 50\,\% do ótimo~\cite{choi2009}. Com grafo completo (todos falam com todos), o
resultado coincide com o do \emph{guloso global}: fixar o par de maior lance entre todos, remover o interceptador e a pista,
e repetir.
\adap Implementamos as duas formas: o guloso global centralizado (M10) e o CBAA iterativo com grafo completo (M11),
representando a comunicação via estação central de~\cite{feitosa2025}; a Seção~\ref{sec:validacao} confirma que dão a mesma
solução. M12 restringe a comunicação a interceptadores a até 3~km uns dos outros e a 3 rodadas de consenso por decisão.
Com isso, interceptadores de bases diferentes podem acreditar que ganharam a mesma pista; o simulador resolve o conflito
pela ordem de chegada, e o perdedor espera a próxima rodada --- que é o custo real de uma visão inconsistente.
\parm M11: até 100 rodadas de consenso; M12: alcance de 3~km e 3 rodadas.
\espera O guloso global é melhor que os gulosos individuais (não depende da ordem da fila), mas pode perder para o húngaro
(Tabela~\ref{tab:exemplo}). M12 mede quanto se perde com comunicação limitada.
\end{metodo}

\begin{algorithm}[H]
\caption{CBAA: consenso sobre lances, uma tarefa por agente (M11, M12).}\label{alg:cbaa}
\begin{algorithmic}[1]
\Require lances $b_{ij}$, grafo $\mathcal{G}$, máximo de rodadas $R$
\State $x_i\gets\varnothing$, $y_{ij}\gets-\infty$, $z_{ij}\gets\varnothing$ para todo $i,j$
\For{$r=1,\dots,R$}
  \For{cada $i$ com $x_i=\varnothing$} \Comment{lances}
    \State $H\gets\{j: b_{ij}>y_{ij}\}$; se $H\neq\varnothing$: $x_i\gets\argmax_{j\in H}b_{ij}$, $y_{ix_i}\gets b_{ix_i}$, $z_{ix_i}\gets i$
  \EndFor
  \For{cada $i$} \Comment{consenso}
    \State $y_{ij},z_{ij}\gets$ maior lance e vencedor entre $\mathcal{N}_i\cup\{i\}$, para todo $j$
    \State se $z_{ix_i}\ne i$: $x_i\gets\varnothing$
  \EndFor
  \State se nada mudou, pare
\EndFor
\Ensure tarefa $x_i$ de cada agente
\end{algorithmic}
\end{algorithm}

% ---------------------------------------------------------------------------------------------------------------
\subsection{WTA: retorno marginal, solução exata e busca local (M13--M16)}
Os métodos desta família otimizam o valor esperado destruído~\eqref{eq:F}, e não o tempo. Isso traz duas diferenças
práticas em relação aos anteriores: eles podem decidir \emph{não lançar}, quando nenhum par compensa o custo de
lançamento; e entendem salvas, isto é, sabem que o segundo interceptador contra a mesma pista vale menos que o primeiro,
porque só conta se o primeiro falhar.

\begin{metodo}{M13 -- Máximo retorno marginal (MMR)}
\ideia Designar um interceptador de cada vez ao par que mais aumenta o valor esperado destruído, considerando o que já
foi designado.
\orig den Broeder, Ellison e Emerling~\cite{denbroeder1959}.
\func O \emph{ganho marginal} de designar $i$ à pista $j$ é quanto o valor esperado destruído aumenta com essa designação,
descontado o custo de lançamento:
\begin{equation}
  \Delta_{ij}=V_j\,\tilde S_j\,q_{ij}-\ell_i,
  \label{eq:mmr}
\end{equation}
em que $\tilde S_j$ é a chance de a pista ainda sobreviver dadas as designações feitas até agora. O algoritmo escolhe o par
livre de maior ganho, designa-o e atualiza a sobrevivência da pista, $\tilde S_j\leftarrow\tilde S_j(1-q_{ij})$ --- de modo
que um segundo interceptador na mesma pista terá ganho menor ---, e repete até que nenhum ganho seja positivo
(Algoritmo~\ref{alg:mmr}). Esse critério de parada é o que implementa a decisão de esperar.
\prop Ótimo quando todas as armas são idênticas; em geral, uma heurística de bom desempenho e ponto de partida de métodos
melhores~\cite{ahuja2007}. Custo de $O(n^2K)$.
\adap Valor e probabilidade de sucesso dados por (E3) e (E4).
\parm $P_T$, $\tau_q$ e $\ell_i$ de (E4).
\espera A qualidade depende do modelo de sucesso~\eqref{eq:q}: se as manobras reduzem a chance real de contato de modo
diferente do previsto pelo modelo, a ordenação dos ganhos fica errada.
\end{metodo}

\begin{algorithm}[H]
\caption{Máximo retorno marginal (M13).}\label{alg:mmr}
\begin{algorithmic}[1]
\Require $q_{ij}$, $V_j$, $S_j$, $\ell_i$, $n_j^0$, $\bar n$
\State $\tilde S_j\gets S_j$, $n_j\gets n_j^0$, todos os $i$ livres
\Loop
  \State $(i^\star,j^\star)\gets\argmax\{V_j\tilde S_jq_{ij}-\ell_i:\ i\text{ livre},\ n_j<\bar n\}$
  \State se o ganho é $\le0$, pare
  \State $x_{i^\star j^\star}\gets1$;\quad $\tilde S_{j^\star}\gets\tilde S_{j^\star}(1-q_{i^\star j^\star})$;\quad $n_{j^\star}\gets n_{j^\star}+1$
\EndLoop
\Ensure $\bm{x}$
\end{algorithmic}
\end{algorithm}

\begin{metodo}{M14, M15 -- WTA exato (húngaro e MILP)}
\ideia Quando cada pista recebe no máximo um interceptador, o objetivo de valor esperado se torna linear e pode ser
otimizado exatamente.
\orig Formulação exata do WTA~\cite{ahuja2007}; M15 representa a família de otimização inteira mista
de~\cite{afonso2020,cruz2025}.
\func Com $\bar n=1$, cada pista tem um único fator no produto de~\eqref{eq:F}, e o valor destruído por designar $i$ a $j$
passa a ser simplesmente $V_jS_jq_{ij}$. O problema fica
\begin{equation}
  \max_{x}\sum_{i,j}\big(V_jS_jq_{ij}-\ell_i\big)x_{ij}\quad\text{s.a.}\quad\sum_jx_{ij}\le1,\ \ \sum_ix_{ij}\le\bar n-n_j^0,\ \
  x_{ij}\in\{0,1\},
  \label{eq:wta11}
\end{equation}
que é de novo um problema de atribuição, agora maximizando valor em vez de minimizar tempo. Resolvemo-lo de duas formas.
M14 usa o húngaro: replica a coluna de cada pista conforme a sua capacidade restante e acrescenta, para cada
interceptador, uma coluna ``esperar'' de valor nulo, de modo que esperar seja escolhido quando nenhum par tem valor
positivo. M15 escreve~\eqref{eq:wta11} como \emph{programa linear inteiro} --- um problema de otimização linear em que as
variáveis devem ser inteiras --- e o resolve com o solver aberto HiGHS, que usa \emph{ramificação e corte}: resolve
relaxações contínuas, divide o problema em ramos quando a solução não é inteira e acrescenta restrições (cortes) que
eliminam soluções fracionárias.
\prop Ambos encontram o ótimo de~\eqref{eq:wta11}; devem atingir o mesmo valor, diferindo só no custo computacional e,
quando há empates, em qual das soluções ótimas escolhem.
\adap Em~\cite{cruz2025}, a alocação está acoplada ao planejamento da trajetória, com obstáculos e bateria; aqui reproduzimos
só a parte de alocação, com a previsão de velocidade constante embutida em $c_{ij}$.
\parm M15: limite de 5~s por rodada.
\espera Diferem do húngaro ponderado (M7) no critério: valor esperado em vez de tempo, e possibilidade de esperar.
\end{metodo}

\begin{metodo}{M16 -- Busca local (VLSN simplificada)}
\ideia Partir de uma boa solução gulosa e melhorá-la por pequenas modificações enquanto houver ganho.
\orig Busca em vizinhança de grande escala de Ahuja et al.~\cite{ahuja2007}.
\func Uma \emph{vizinhança} é o conjunto de soluções obtidas a partir da atual por uma modificação simples. Usamos duas:
\emph{realocação} (um interceptador muda para outra pista viável ou passa a esperar) e \emph{troca} (dois interceptadores
trocam de pista entre si). A partir da solução do MMR, avaliamos todas as soluções das duas vizinhanças, adotamos a de
maior valor de~\eqref{eq:F} se ela for melhor que a atual, e repetimos até não haver melhoria.
\prop Termina num \emph{ótimo local}: uma solução que nenhuma realocação ou troca isolada melhora. Nunca é pior que o MMR.
\adap A VLSN original usa vizinhanças exponencialmente grandes (reatribuições em cadeia) exploradas por algoritmos de fluxo
em rede; a nossa usa vizinhanças de tamanho polinomial, avaliadas de forma vetorizada.
\parm Até 30 iterações.
\end{metodo}

% ---------------------------------------------------------------------------------------------------------------
\subsection{Metaheurísticas (M17--M20)}
As quatro metaheurísticas otimizam o mesmo objetivo~\eqref{eq:F}. Todas representam uma solução por um vetor
$\bm{g}\in\{-1,0,\dots,K-1\}^{|\mathcal{D}_t|}$, com uma posição por interceptador da fila contendo o índice da pista
designada ($-1$ significa esperar), e só consideram pares viáveis. Todas incluem a solução do MMR no ponto de partida, de
modo que nunca terminam abaixo dela. Por serem estocásticas, usam sementes aleatórias fixadas pelo instante da rodada,
para que os resultados sejam reprodutíveis.

\begin{metodo}{M17 -- Algoritmo genético}
\ideia Imitar a seleção natural: manter uma \emph{população} de soluções, combinar as melhores para gerar novas e
introduzir variações aleatórias.
\orig Lee, Su e Lee~\cite{lee2003ga}, que semeiam a população com soluções gulosas (``eugenia gulosa'').
\func A cada \emph{geração}: (i) \emph{seleção por torneio} --- sorteiam-se três soluções e a melhor delas é escolhida
como pai, o que favorece as boas sem descartar as demais; (ii) \emph{cruzamento uniforme} --- o filho herda a designação de
cada interceptador de um dos dois pais, sorteado com igual chance; (iii) \emph{mutação} --- cada designação do filho é
trocada por uma opção viável sorteada com pequena probabilidade, o que mantém a diversidade; (iv) \emph{elitismo} --- as
duas melhores soluções passam intactas para a próxima geração, para não perder o melhor já encontrado. A qualidade
(\emph{aptidão}) de cada solução é~\eqref{eq:F}.
\parm População de 40, 40 gerações, torneio de 3, probabilidade de mutação $1/|\mathcal{D}_t|$, elitismo de 2.
\end{metodo}

\begin{metodo}{M18 -- Colônia de formigas}
\ideia Imitar formigas que, ao encontrar comida, reforçam com feromônio o caminho percorrido; caminhos bons acumulam mais
feromônio e atraem mais formigas.
\orig Lee, Lee e Su~\cite{lee2002aco}.
\func Cada \emph{formiga} constrói uma solução visitando os interceptadores em ordem aleatória e escolhendo, para cada um,
uma opção $j$ (uma pista ou esperar) com probabilidade
\begin{equation}
  P_{ij}=\frac{\tau_{ij}^{\alpha}\eta_{ij}^{\beta}}{\sum_{k}\tau_{ik}^{\alpha}\eta_{ik}^{\beta}},
\end{equation}
que combina duas informações: a \emph{atratividade heurística} $\eta_{ij}=\max(V_jS_jq_{ij}-\ell_i,10^{-6})$, o ganho do
par se ele fosse o único, e o \emph{feromônio} $\tau_{ij}$, que acumula a experiência das iterações anteriores. Os
expoentes $\alpha$ e $\beta$ regulam o peso de cada uma. Pistas já no limite de capacidade recebem probabilidade nula. Ao
fim de cada iteração, todo o feromônio \emph{evapora} parcialmente (o que permite esquecer escolhas ruins) e a melhor
solução encontrada é reforçada: $\tau\leftarrow(1-\rho)\tau+\rho F^\star\mathbf{1}_{\bm{g}^\star}$, em que
$\mathbf{1}_{\bm{g}^\star}$ marca os pares da melhor solução e $F^\star$ é o seu valor.
\adap A opção ``esperar'' é tratada como um caminho com atratividade fixa ($\eta_{i,-1}=0{,}05$).
\parm 16 formigas, 15 iterações, $\alpha=1$, $\beta=2$, taxa de evaporação $\rho=0{,}1$.
\end{metodo}

\begin{metodo}{M19 -- Recozimento simulado}
\ideia Imitar o resfriamento lento de um metal: no início, o sistema ``quente'' aceita mudanças que pioram a solução, o que
permite escapar de ótimos locais; à medida que esfria, torna-se cada vez mais exigente.
\orig Metaheurística de uso amplo no WTA~\cite{kline2019}.
\func Partindo do MMR, propõe-se a cada passo uma realocação aleatória de um interceptador. Uma melhoria é sempre aceita;
uma piora $\Delta F<0$ é aceita com probabilidade $e^{\Delta F/T}$, em que $T$ é a \emph{temperatura}: com $T$ alta, quase
toda piora é aceita; com $T$ baixa, quase nenhuma. A temperatura diminui a cada passo, e guarda-se a melhor solução
visitada.
\parm 1500 movimentos, temperatura inicial $T_0=0{,}2|F_0|$ (proporcional ao valor da solução inicial) e resfriamento
$T\leftarrow0{,}995\,T$.
\end{metodo}

\begin{metodo}{M20 -- Enxame de partículas}
\ideia Imitar um bando: cada solução candidata (uma \emph{partícula}) se move pelo espaço de soluções, atraída pela melhor
posição que ela própria já visitou e pela melhor posição encontrada pelo bando.
\orig Metaheurística de uso amplo no WTA~\cite{kline2019}.
\func Cada partícula tem uma posição $\bm{X}$ e uma velocidade $\bm{V}$, atualizadas por
\begin{equation}
  \bm{V}\leftarrow\omega\bm{V}+c_1\bm{r}_1\odot(\bm{X}^{p}-\bm{X})+c_2\bm{r}_2\odot(\bm{X}^{g}-\bm{X}),\qquad
  \bm{X}\leftarrow\mathrm{clip}_{[0,1]}(\bm{X}+\bm{V}),
\end{equation}
em que $\bm{X}^p$ é a melhor posição da própria partícula, $\bm{X}^g$ a melhor do bando, $\bm{r}_1,\bm{r}_2$ números
aleatórios uniformes e $\odot$ o produto elemento a elemento. O primeiro termo é a \emph{inércia} (manter a direção), o
segundo a atração \emph{individual} e o terceiro a atração \emph{social}.
\adap O enxame de partículas trabalha com posições contínuas, mas a alocação é discreta. Usamos a codificação por
\emph{chaves aleatórias}: cada partícula é uma tabela de números em $[0,1]$, com uma linha por interceptador e uma coluna
por opção (cada pista e esperar); para obter a alocação, cada interceptador recebe a opção viável de maior número na sua
linha. Uma partícula começa na codificação da solução do MMR.
\parm 20 partículas, 25 iterações, inércia $\omega=0{,}7$, $c_1=c_2=1{,}5$.
\end{metodo}

\paragraph{O que esperar das metaheurísticas.} Elas só podem superar o MMR se o objetivo~\eqref{eq:F} tiver estrutura que
a escolha gulosa perde; em troca, custam uma a duas ordens de grandeza mais tempo por rodada
(Seção~\ref{sec:validacao}).

% ---------------------------------------------------------------------------------------------------------------
\subsection{Alocação dinâmica (M21 e variantes)}

\begin{metodo}{M21 -- Húngaro com reserva}
\ideia Quando as ameaças chegam em ondas, não gastar todos os interceptadores na primeira.
\orig Lição do WTA dinâmico~\cite{hosein1990}.
\func Resolve o problema de atribuição~\eqref{eq:lap} como M7, mas aplica uma regra de reserva: um interceptador da base só
é lançado contra uma pista de baixa prioridade ($w_j<w_P$) se, depois do lançamento, restarem mais de 20\,\% dos
interceptadores ociosos. Pistas dentro de zonas prioritárias podem usar a reserva.
\adap É uma política de horizonte deslizante: aproxima o valor de guardar recursos sem modelar explicitamente as ondas
futuras, que o defensor não conhece.
\parm Reserva de 20\,\% da frota.
\espera Pode perder ameaças de baixa prioridade na primeira onda para ganhar nas seguintes.
\end{metodo}

\paragraph{\emph{Shoot--look--shoot} e salvas.} Todos os métodos permitem trocar de alvo em voo: um interceptador que perde
a pista ou percebe que ela era falsa é redesignado em vez de desperdiçado. Isso é a forma \emph{shoot--look--shoot} do WTA
dinâmico~\cite{hosein1990}: observar antes de reengajar. A alternativa, \emph{shoot--shoot}, manda dois interceptadores à
mesma pista de uma vez ($\bar n=2$), trocando interceptadores por chance de sucesso. Medimos as duas por ablação
(Seção~\ref{sec:res_ablacao}): desligando a troca de alvo, e permitindo salvas no M13 e no M17.

% ---------------------------------------------------------------------------------------------------------------
\subsection{Arquiteturas comportamentais (M22, M23)}
Nas árvores de comportamento, cada interceptador executa a sua própria árvore de forma independente, com o quadro de radar
comum; quando dois escolhem a mesma pista, prevalece o primeiro numa ordem fixa. Diferentemente dos métodos conjuntos, a
árvore decide \emph{se} lança, além de \emph{contra quem} --- o que lhe permite poupar interceptadores.

\begin{metodo}{M22 -- Árvore com alocação por distância}
\ideia Engajar só ameaças que já estão perto, e sempre a mais próxima.
\orig Ricardo Jr et al.~\cite{ricardo2023}.
\func A árvore tem um seletor com dois ramos, tentados em ordem: \emph{engajar}, se existe ameaça livre e plausível a menos
de $R_e$ do interceptador (o alvo é a mais próxima); senão, \emph{esperar}. Em voo, o interceptador mantém o alvo até
perdê-lo.
\adap No original, os defensores protegem um drone líder numa arquitetura líder--seguidor; aqui cada interceptador decide
sozinho a partir da sua base.
\parm Raio de engajamento $R_e=4$~km.
\espera Economiza interceptadores, mas, ao esperar que a ameaça se aproxime, deixa pouco tempo para alcançar alvos rápidos.
\end{metodo}

\begin{metodo}{M23 -- Árvore com utilidade afim}
\ideia Decidir se vale a pena avançar sobre uma ameaça por uma nota de \emph{utilidade} que pesa o custo de ir, o perigo
da ameaça e a pressão total sobre a defesa.
\orig Feitosa, Ricardo Jr e Santos~\cite{feitosa2025}.
\func A árvore original tem três ramos: \emph{responder} (inspecionar a ameaça e então avançar ou pedir reforço),
\emph{reforço} (atender pedidos de apoio de outros defensores) e \emph{ocioso}. A decisão de avançar é dada por uma
utilidade afim --- uma soma ponderada de fatores mais uma constante. Na nossa versão, o interceptador toma como alvo a
ameaça livre de menor custo $c_{ij}$ e avança se
\begin{equation}
  U_i = c_1\frac{c_{ij}}{100} + c_2\,\theta'_j + c_3\frac{\Theta}{N} + c_4 > 0,\qquad
  \theta'_j=\frac{w_j}{w_P}e^{-t^Z_j/120},\quad \Theta=\sum_{k}\theta'_k .
  \label{eq:util}
\end{equation}
Os quatro termos correspondem aos da utilidade original: o tempo de deslocamento (normalizado por 100~s), a ameaça do alvo
$\theta'_j$ (que vale até 1 e decai com o tempo até a zona, em escala de 2~minutos), a ameaça composta $\Theta$ (a soma das
ameaças de todas as pistas, dividida pelo número de aeronaves, que mede a pressão sobre a defesa) e uma constante que
representa a aversão a risco. Os coeficientes $c_1,\dots,c_4$ definem o comportamento: um $c_1$ negativo desestimula
perseguir alvos distantes; $c_2$ e $c_3$ positivos estimulam reagir a ameaças fortes ou a muita pressão.
\adap O ramo de reforço trata alvos que resistem a vários disparos e não se aplica a interceptadores que neutralizam com um
único contato. A ameaça de uma pista usa o tempo até a zona prioritária. Como no original, os coeficientes são ajustados
por um algoritmo genético; aqui, para minimizar o dano~\eqref{eq:dano} médio em cenários de treino.
\parm Algoritmo genético com população de 12, 8 gerações, torneio, cruzamento aritmético (média ponderada dos pais) e
mutação gaussiana de 10\,\% da faixa de cada coeficiente, avaliado em 10 cenários de treino.
\end{metodo}

% ---------------------------------------------------------------------------------------------------------------
\subsection{Políticas aprendidas (M24--M29)}\label{sec:redes}

\paragraph{Ideia geral.} Uma política aprendida é uma rede neural que recebe a situação vista por um interceptador e
devolve uma distribuição de probabilidade sobre as ações possíveis. Em princípio, ela pode capturar regularidades que as
regras ignoram --- por exemplo, que certas pistas quase sempre são falsas, ou que vale esperar uma pista se confirmar antes
de lançar. Todos os interceptadores consultam a mesma rede, cada um com a sua observação (\emph{política compartilhada}).

\paragraph{Observação e ação.} Para o interceptador da vez, a rede recebe até $K=64$ pistas livres. Cada pista é descrita
por 17 \emph{traços de decisão} --- posição e velocidade relativas ao interceptador, tempo previsto de interceptação, se é
alcançável com a autonomia restante, idade, incerteza, se já está coberta por outro interceptador, altura, razão de
subida, distância e velocidade de aproximação à zona prioritária, se é o alvo atual --- e por 10 \emph{traços de
legitimidade}, só de radar, que ajudam a separar aeronaves de falsos alarmes: velocidade, razão de subida, altura,
incerteza, idade, tempo de vida, número e taxa de medidas e distância ao radar mais próximo. Há ainda 7 \emph{traços
globais}: quantos interceptadores estão pendentes e livres, quantas pistas há, o instante, a autonomia restante, se o
interceptador está em voo e a idade da pista que persegue. A ação é uma das pistas ou a ação nula.

\paragraph{Classificador de legitimidade.} Em vez do filtro (E1), a rede tem um pequeno classificador $\lambda$ que
recebe os traços de legitimidade de uma pista e devolve a probabilidade de ela ser uma aeronave. Ele é treinado com o
rótulo verdadeiro, que o simulador conhece (mas que nunca é dado à rede na decisão), por \emph{entropia cruzada binária}
--- a perda padrão de classificação, que penaliza atribuir probabilidade baixa à classe correta ---, com as classes
balanceadas, porque os falsos alarmes são maioria. A probabilidade estimada entra na nota de cada pista.

\begin{metodo}{Arquitetura DeepSets}
\orig Zaheer et al.~\cite{zaheer2017}.
\func O número de pistas varia e a ordem em que são listadas é arbitrária; a decisão não pode depender dela. A
arquitetura DeepSets resolve isso processando cada pista com a mesma rede $\phi$ (um \emph{perceptron multicamadas}, MLP,
a rede neural mais simples) e combinando os resultados por uma média, operação que não depende da ordem:
\begin{equation}
  \bm{h}_j=\phi([\bm{o}_j,\bm{l}_j,\bm{g}]),\quad \bar{\bm{h}}=\tfrac{1}{|\mathcal{V}|}\textstyle\sum_{j}\bm{h}_j,\quad
  z_j=\psi\big([\bm{h}_j,\bar{\bm{h}},\sigma(\lambda(\bm{l}_j))]\big),\quad z_{\varnothing}=\chi([\bar{\bm{h}},\bm{g}]),
\end{equation}
em que $\bm{o}_j$, $\bm{l}_j$ e $\bm{g}$ são os traços de decisão, de legitimidade e globais; $\bm{h}_j$ é o
\emph{código} da pista; $\bar{\bm{h}}$, o \emph{contexto}, resume todas as pistas; $z_j$ é a nota da pista, que depende do
seu código, do contexto e da probabilidade de ser legítima; e $z_\varnothing$ é a nota da ação nula. As notas são
convertidas em probabilidades pela função \emph{softmax}, $\pi_j=e^{z_j}/\sum_ke^{z_k}$.
\parm MLPs de 96 unidades; classificador com camadas de 10, 64, 64 e 1 unidades.
\end{metodo}

\begin{metodo}{Arquitetura com atenção}
\orig Vaswani et al.~\cite{vaswani2017}; aplicação a problemas combinatórios por Kool et al.~\cite{kool2019}; agregação por
atenção em grafos em Jia et al.~\cite{jia2026}.
\func Na média do DeepSets, todas as pistas pesam igual no contexto. Com \emph{atenção}, cada pista ``consulta'' as demais
e decide quanto cada uma importa para ela. Cada código $\bm{h}_j$ gera uma consulta $\bm{q}_j$, uma chave $\bm{k}_j$ e um
valor $\bm{v}_j$ (projeções lineares aprendidas); o peso que a pista $j$ dá à pista $k$ é proporcional a
$e^{\bm{q}_j\cdot\bm{k}_k/\sqrt{d}}$, e o novo código de $j$ é a média dos valores ponderada por esses pesos:
\begin{equation}
  \text{Atenção}(Q,K,V)=\mathrm{softmax}\!\Big(\frac{QK^\top}{\sqrt{d}}\Big)V .
\end{equation}
Assim, o código de uma pista passa a refletir, por exemplo, se há outra pista mais fácil perto dela. Usamos duas camadas
de \emph{codificador Transformer} (atenção seguida de um MLP, com conexões residuais e normalização), com 4
\emph{cabeças} de atenção em paralelo, antes da média; o restante é igual ao DeepSets. A operação continua independente da
ordem das pistas.
\parm 2 camadas, 4 cabeças, dimensão 96.
\end{metodo}

\begin{metodo}{M24, M25 -- Imitação (DAgger)}
\ideia Aprender a decidir como um especialista, mas mais rápido, a partir de exemplos rotulados por ele.
\orig Ross, Gordon e Bagnell~\cite{ross2011}.
\func O especialista é o húngaro ponderado (M7). Uma imitação ingênua treinaria a rede só com situações produzidas pelo
especialista; mas, ao ser usada, a rede comete pequenos erros que a levam a situações que o especialista nunca produziu, e
nessas ela erra mais --- um efeito chamado \emph{mudança de distribuição}. O DAgger corrige isso iterativamente: na primeira
iteração, o especialista conduz os episódios e rotula as decisões; nas seguintes, a \emph{rede} conduz, e o especialista
apenas rotula a decisão certa em cada situação que ela visitou. Os exemplos de todas as iterações são acumulados, e a rede
é retreinada para prever a ação do especialista, minimizando a entropia cruzada entre a sua distribuição e a ação rotulada.
\adap Quase todas as decisões do especialista são ``esperar'', porque a cada rodada todos os interceptadores ociosos nas
bases decidem e a maioria espera. Treinada com essa proporção, uma versão preliminar da rede aprendeu a esperar sempre.
Guardamos então todas as decisões de lançar ou trocar e só 10\,\% das de esperar, sorteadas. Para não enviesar a rede no
sentido oposto, cada exemplo de esperar recebe \emph{peso de importância} 10 (o inverso da fração guardada), de modo que a
perda média é a mesma que se teria com todos os exemplos:
\begin{equation}
  \mathcal{L}_{\mathrm{IL}}=\frac{\sum_k\omega_k\,\mathrm{CE}\big(\pi(\cdot\mid o_k),a^\star_k\big)}{\sum_k\omega_k}
  +0{,}5\,\mathcal{L}_{\mathrm{leg}},\qquad\omega_k=\begin{cases}10,&a^\star_k=\text{esperar},\\1,&\text{caso contrário,}\end{cases}
\end{equation}
em que $\mathrm{CE}$ é a entropia cruzada, $a^\star_k$ a ação do especialista no exemplo $k$ e $\mathcal{L}_{\mathrm{leg}}$ a
perda do classificador de legitimidade.
\parm 5 iterações de 28 episódios, 8 épocas de treino por iteração, otimizador Adam com taxa $10^{-3}$, 3 sementes de
treino.
\espera A rede pode, no melhor caso, igualar o especialista; a vantagem é a velocidade de decisão, e o interesse é servir
de ponto de partida para o reforço.
\end{metodo}

\begin{metodo}{M26--M29 -- Reforço (PPO)}
\ideia Aprender diretamente a minimizar o dano, por tentativa e erro, sem depender de um especialista.
\orig Otimização de política proximal (PPO)~\cite{schulman2017}; políticas centralizadas para alocação por
reforço~\cite{liu2023drones,palmas2025,jia2026}.
\func A rede joga episódios completos e recebe, a cada passo, uma \emph{recompensa} que premia a interceptação de
ameaças armadas, pune impactos pelo seu peso, cobra cada interceptador desperdiçado (inclusive os gastos em iscas) e
cobra um pequeno custo pelo tempo com ameaças ainda vivas, para estimular rapidez:
\begin{equation}
  r_t = \sum_{a:\,\nu_a>0}\mathbf{1}\{a\text{ interceptada}\} - \sum_a w_a\mathbf{1}\{a\text{ atinge o impacto}\}
  - \eta\,\Delta W_t - \frac{N_{\mathrm{vivas}}(t)\,\Delta t}{300\,N},
\end{equation}
em que $\Delta W_t$ é o número de interceptadores desperdiçados no passo. Os dois primeiros termos e o terceiro
correspondem a~\eqref{eq:dano} e~\eqref{eq:custo_ep}: a recompensa ensina a proteger sem disparar à toa.
O objetivo é maximizar a soma das recompensas. Para isso, além da política, a rede tem uma cabeça de \emph{valor}, que
estima quanto de recompensa ainda se espera obter a partir de uma situação. A \emph{vantagem} $\hat A_t$ de uma ação é
quanto ela foi melhor que o esperado: positiva se a recompensa obtida superou a estimativa de valor. Usamos a estimativa
GAE, uma média ponderada das diferenças entre recompensa observada e valor previsto ao longo dos passos seguintes, que
equilibra ruído e viés. A atualização torna mais prováveis as ações de vantagem positiva e menos as de vantagem negativa.
O PPO limita o tamanho de cada atualização para não destruir o que já foi aprendido:
\begin{equation}
  \mathcal{L}^{\mathrm{CLIP}}=\mathbb{E}_t\Big[\min\big(\rho_t\hat A_t,\ \mathrm{clip}(\rho_t,1-\epsilon,1+\epsilon)\hat A_t\big)\Big],
  \qquad\rho_t=\frac{\pi_{\mathrm{nova}}(a_t\mid o_t)}{\pi_{\mathrm{antiga}}(a_t\mid o_t)},
\end{equation}
em que $\rho_t$ mede quanto a nova política mudou a probabilidade da ação tomada; o recorte ($\mathrm{clip}$) ignora
mudanças maiores que $\pm\epsilon$. Somam-se uma perda de valor, um bônus de \emph{entropia} (que estimula a política a
continuar explorando ações diferentes) e a perda de legitimidade.
\adap M26 e M27 partem das redes de imitação (M24, M25), com taxa de aprendizado menor, para refinar o comportamento já
aprendido; M28 e M29 partem do zero. A perda de legitimidade é mantida como tarefa auxiliar.
\parm 40 atualizações de 7 episódios cada, 4 passadas pelos dados por atualização, lotes de 512 exemplos, $\epsilon=0{,}2$,
fator de desconto $\gamma=0{,}995$, parâmetro do GAE 0,95, bônus de entropia 0,01; taxa de aprendizado $10^{-4}$ (M26, M27)
ou $3\times10^{-4}$ (M28, M29); 3 sementes de treino, só nos cenários de treino.
\espera O reforço do zero é o mais custoso de treinar e o mais variável entre sementes; o refino a partir da imitação
tende a ser mais estável.
\end{metodo}

% ---------------------------------------------------------------------------------------------------------------
\subsection{Resumo de custo e garantias}
A Tabela~\ref{tab:metodos} resume, para cada família, o objetivo otimizado, a garantia conhecida e o custo por rodada.

\begin{table}[t]
\centering
\caption{Objetivo, garantia e custo por rodada. $n$: interceptadores da rodada; $K$: pistas; $G$: iterações; $P$: população;
$R$: rodadas de consenso.}
\label{tab:metodos}
\small
\begin{tabularx}{\textwidth}{@{}l l X l@{}}
\toprule
Métodos & Objetivo & Garantia & Custo por rodada \\
\midrule
M2--M4 & tempo (individual) & nenhuma; depende da ordem da fila & $O(nK)$ \\
M5 & ameaça, depois tempo & nenhuma & $O(K\log K+nK)$ \\
M6--M8 & tempo total \eqref{eq:lap} & ótimo da rodada & $O(n^2K)$ \\
M9 & tempo total \eqref{eq:lap} & a $n\varepsilon$ do ótimo & pseudo-polinomial \\
M10--M11 & lances $w_j/c_{ij}$ & $\ge50\,\%$ do ótimo~\cite{choi2009} & $O(R\,n^2K)$ \\
M12 & idem, comunicação local & sem garantia de ausência de conflito & $O(3\,n^2K)$ \\
M13 & valor esperado \eqref{eq:F} & ótimo para armas idênticas & $O(n^2K)$ \\
M14, M15 & valor esperado, $\bar n=1$ & ótimo da rodada & $O(n^2K)$ / exponencial no pior caso \\
M16 & valor esperado & ótimo local, $\ge$ M13 & $O(G(nK+n^2)nK)$ \\
M17--M20 & valor esperado & $\ge$ M13 & $O(GP\,nK)$ \\
M21 & tempo + reserva & nenhuma & $O(n^2K)$ \\
M22, M23 & regra / utilidade & nenhuma & $O(K)$ por interceptador \\
M24--M29 & imitação / recompensa & nenhuma & uma inferência por interceptador \\
\bottomrule
\end{tabularx}
\end{table}
